\bookchapter{第二章：整理你嘅數碼房間——檔案管理}

\chapterintro{手機用耐咗會亂，好習慣由管理檔案開始。就好似你嘅房間需要定期整理一樣，你嘅數碼空間都需要打理。一個井井有條嘅手機，唔只揾嘢方便，運行都更流暢。}

\booksection{一、三層存儲體系：你嘅數碼倉庫}

\begin{enumerate}
  \item \textbf{機身存儲}：核心區，速度快，系統同App都喺度。呢個係你手機最寶貴嘅空間，要精打細算。
  \item \textbf{SD卡（如果支援）}：擴展衣櫃，速度慢啲，適合放電影、相片等大檔案。注意：而家好多旗艦機唔支援SD卡喇。
  \item \textbf{雲存儲}：保險櫃+同步網盤。防遺失，都可以喺多個設備間同步。iCloud、百度網盤、OneDrive都係常用選擇。
\end{enumerate}

\booksection{二、認識「檔案管理器」}

安卓一般有個 [檔案] 或 [檔案管理] App，iOS係自帶嘅 [檔案] App。佢哋係你查看手機所有檔案嘅「資源管理器」。

\booksubsection{常見資料夾含義：}

\begin{itemize}
  \item \textbf{DCIM}：相機拍攝嘅相片同影片。
  \item \textbf{Pictures}：其他圖片，如截圖、儲存嘅圖片。
  \item \textbf{Download}：瀏覽器下載嘅檔案。
  \item \textbf{Music}：音樂檔案。
  \item \textbf{Movies}：影片檔案。
  \item \textbf{Documents}：文檔檔案。
  \item \textbf{Android}：安卓系統資料夾，入面係各App嘅數據，一般唔好手動刪除。
\end{itemize}

\booksection{三、檔案系統結構：深入瞭解}

安卓基於Linux，檔案系統有一定嘅層級結構。雖然普通用戶唔需要深入瞭解，但知道一啲基本概念有幫助。

\booksubsection{存儲分區：}

\begin{itemize}
  \item \textbf{/system}：系統分區，存放作業系統檔案。普通用戶無權訪問，需要root權限。
  \item \textbf{/data}：數據分區，存放用戶數據同應用數據。呢個係你日常使用嘅空間。
  \item \textbf{/sdcard 或 /storage/emulated/0}：用戶可訪問嘅存儲空間，你嘅相片、下載檔案都喺度。
\end{itemize}

\booksubsection{App數據存儲：}

\begin{itemize}
  \item \textbf{內部存儲}：App嘅私有空間，其他App無法訪問。卸載App嗰陣會被刪除。
  \item \textbf{外部存儲}：App喺用戶可訪問區域存儲嘅檔案，如相片、下載內容。卸載App嗰陣唔一定會被刪除。
\end{itemize}

\booksection{四、數據備份與恢復：畀你嘅數據上保險}

數據無價，備份係必須養成嘅習慣。

\booksubsection{備份策略：}

\begin{enumerate}
  \item \textbf{雲備份}：自動、方便，但需要網絡同雲空間。iCloud、華為雲、小米雲等都支援。
  \item \textbf{本地備份}：用數據線連接電腦，用官方工具備份。速度快、空間大，但需要手動操作。
  \item \textbf{雙重備份}：雲備份+本地備份，最安全。
\end{enumerate}

\booksubsection{備份乜嘢？}

\begin{itemize}
  \item \textbf{相片同影片}：最佔空間，都最珍貴。
  \item \textbf{通訊錄}：存在雲帳號入面最方便。
  \item \textbf{微信聊天記錄}：用微信自帶嘅備份功能。
  \item \textbf{重要文檔}：存在雲盤或電腦上。
\end{itemize}

\booksection{五、檔案安全：保護你嘅私隱}

\begin{itemize}
  \item \textbf{敏感檔案加密}：好多手機有「安全資料夾」或「私密空間」功能，可以將敏感檔案放入去，需要密碼或指紋先可以訪問。
  \item \textbf{刪除唔等於銷毀}：普通刪除淨係標記空間可覆蓋，數據仲喺度。要徹底刪除，需要用專門嘅安全擦除工具，或者多次覆蓋寫入。
  \item \textbf{雲端安全}：選擇信譽好嘅雲服務商，開啟雙重認證，定期檢查登錄設備。
\end{itemize}

\booksection{六、空間唔夠？嚟場大掃除}

\begin{itemize}
  \item \textbf{清理快取}：App嘅臨時記憶，清咗可以騰出地方。喺設定 > 存儲 入面可以清理。
  \item \textbf{查找大檔案}：用檔案管理器按大小排序，揪出佔空間嘅「元兇」。
  \item \textbf{刪除殘留}：卸載App後，有時會留低垃圾檔案，手動清掉。
  \item \textbf{清理相簿}：刪除重複相片、模糊相片、唔需要嘅截圖。
  \item \textbf{清理微信}：微信係空間大戶，定期清理聊天記錄同快取好有必要。
\end{itemize}

\booksection{七、常見檔案格式全解析}

手機入面各種檔案後綴名令人頭大，搞懂佢哋可以幫你更好咁管理檔案：

\begin{center}
\begin{tabularx}{\linewidth}{|X|X|X|X|}
\hline
\textbf{格式} & \textbf{類型} & \textbf{特性} & \textbf{推薦打開方式} \tabularnewline
\hline
JPG/JPEG & 圖片 & 有損壓縮，體積細，兼容性最好。手機影相預設格式。 & 相簿/任意圖片查看器 \tabularnewline
PNG & 圖片 & 無損壓縮，支援透明背景。截圖、圖標常用。體積比JPG大。 & 相簿/任意圖片查看器 \tabularnewline
MP3 & 音頻 & 有損壓縮，體積細，兼容性極好。最通用嘅音頻格式。 & 任意音樂播放器 \tabularnewline
FLAC & 音頻 & 無損壓縮，音質與CD一致，體積比MP3大3-5倍。發燒友首選。 & 音樂播放器/專業播放App \tabularnewline
WAV & 音頻 & 無壓縮原始音頻，音質最好但體積巨大。專業錄音常用。 & 音樂播放器 \tabularnewline
M4A & 音頻 & AAC編碼，比MP3同碼率音質更好。蘋果生態常用。 & Apple Music/音樂播放器 \tabularnewline
MP4 & 影片 & 最通用嘅影片封裝格式，兼容性極好。手機錄製嘅影片預設格式。 & 相簿/任意影片播放器 \tabularnewline
MKV & 影片 & 開源封裝格式，支援多音軌、多字幕。電影資源常用。 & VLC/nPlayer \tabularnewline
DOC/DOCX & 文檔 & Word文檔。DOCX係壓縮格式，體積更細。DOC係舊版格式。 & WPS/Office \tabularnewline
PPT/PPTX & 演示 & PowerPoint演示文稿。PPTX係壓縮格式。 & WPS/Office \tabularnewline
XLS/XLSX & 表格 & Excel電子表格。XLSX係壓縮格式。 & WPS/Office \tabularnewline
PDF & 文檔 & 跨平臺格式，排版固定唔可編輯。合同、論文、電子書常用。 & Adobe Reader/WPS \tabularnewline
ZIP & 壓縮 & 最通用嘅壓縮格式，跨平臺兼容好。 & 系統自帶/解壓App \tabularnewline
RAR & 壓縮 & 壓縮率比ZIP高，但需要專用軟件解壓。 & 解壓App \tabularnewline
ISO & 鏡像 & 光碟鏡像檔案，完整複製光碟內容。系統安裝包常用。 & 解壓App/虛擬光驅 \tabularnewline
HTML & 網頁 & 網頁檔案，用瀏覽器打開。儲存嘅網頁常用呢個格式。 & 瀏覽器 \tabularnewline
\hline
\end{tabularx}
\end{center}

\begin{tcbpractice}{練一練：整理「下載」資料夾}
打開檔案管理器 > 進入「下載」資料夾。睇下入面都有啲乜嘢？將冇用嘅刪掉，有用嘅分門別類（比如新建「工作資料」、「表情包」資料夾放入去）。順便睇下存儲空間，瞭解邊啲類型檔案佔用最多。

\end{tcbpractice}

\booksection{八、圖片與檔案系統管理：等資產井井有條}

上一節我哋學識咗「認識資料夾」，但淨係認識都唔夠。本節要講嘅係\textbf{點樣主動建立一套可長期執行嘅管理規範}，等你嘅檔案隨時揾得到、唔遺失、唔被偷窺。

\booksubsection{檔案分類與命名規範}

混亂嘅根源係冇規則。建立規則其實好簡單，關鍵在於\textbf{堅持}。

\begin{itemize}
  \item \textbf{「時間+事件」命名法}：資料夾名入面帶時間戳同事件描述，方便日後回憶同檢索。例如\texttt{202606\_旅行票據}、\texttt{2025年終總結}、\texttt{2024孩子生日}。咁樣既能睇到時間線，一眼就知道內容。
  \item \textbf{三級分類體系}：
  \begin{itemize}
    \item \textbf{一級}：按檔案類型（文檔 / 圖片 / 影片 / 音頻）
    \item \textbf{二級}：按使用場景（工作 / 生活 / 學習）
    \item \textbf{三級}：按具體事件或項目（如「2026Q2\_項目A」）
  \end{itemize}
  \item \textbf{月度清理計劃}：每月固定一個日期（建議出糧嗰日或月初1號），對「下載」目錄進行集中整理：\textbf{歸檔有用嘅、刪除冇用嘅、轉存珍貴嘅}。一年 12 次小整理，比起攢一年大掃除輕鬆得多。
\end{itemize}

\begin{tcbpractice}{本月任務清單}
揀一個空閒嘅週末，按以下順序操作：

\begin{enumerate}
  \item 喺檔案管理器中創建\texttt{2026Q2\_項目A}形式嘅本月工作資料夾；
  \item 將所有「下載」目錄入面本週嘅檔案過一遍，留低有用嘅、歸檔、刪除冇用嘅；
  \item 畀今年嘅家庭相片建立\texttt{2025家庭}、\texttt{2026家庭}呢種年度資料夾。
\end{enumerate}

\end{tcbpractice}

\booksubsection{冗餘檔案處理策略}

手機用耐咗會「長胖」，多數係冗餘檔案喺度作怪。

\begin{itemize}
  \item \textbf{72 小時自動清理}：微信、QQ 等社交 App 接收嘅臨時檔案（表情包、聊天圖片、短影片），大多睇一次就冇用。喺 [微信] > [設定] > [通用] > [存儲空間] 入面開啟「自動清理」，將過期快取設成 72 小時。手機瞬間清爽唔少，仲唔丟重要記錄。
  \item \textbf{存儲空間定向清理}：安卓 [設定] > [存儲] > [清理加速]，iOS [設定] > [通用] > [iPhone 存儲空間]，系統會按「快取 / 臨時檔案 / 可卸載 App / 大檔案」分類列出。重點清理：(1) 快取檔案；(2) 已過期安裝包（升級 App 後留低嘅舊版安裝包）；(3) 重複檔案（同名同尺寸嘅相片）。
  \item \textbf{自動分類規則}：多數品牌相簿 App 支援智能分類：人物、風景、文檔、截圖。開啟後，新拍嘅合影會自動歸入「人物 > 張三」，掃描嘅合同會自動歸入「文檔」。\textbf{配合每月整理，可極大降低「揾唔到嘢」嘅概率}。
\end{itemize}

\booksubsection{圖片安全處理規範}

圖片係手機入面最「敏感」嘅數據類型——一張隨手拍就可以暴露你嘅位置、設備、時間。我哋建議\textbf{建立「雙重保護」工作流}：

\begin{tcbwarn}{發送證件相/敏感截圖嘅標準流程}
\begin{enumerate}
  \item \textbf{可見資訊打碼}：用系統自帶或第三方工具（如美圖秀秀、Snapseed）\textbf{塗抹身份證號中間幾位、車牌號、門牌、姓名、金額}等關鍵字段；
  \item \textbf{元數據剝離}：發送前喺相簿中點「分享」或「導出」，系統會詢問係咪保留位置資訊，選擇\textbf{「唔保留」}。咁樣拍攝設備型號、GPS 坐標、拍攝時間就不會跟住圖片一齊發出去；
  \item \textbf{臨時水印}（可選）：喺敏感截圖上加「僅供XX辦理使用」嘅文字水印，避免被濫用。
\end{enumerate}

\end{tcbwarn}

元數據（EXIF）清理嘅常用方法（基礎操作，深度防護見第六章「相片脱敏處理」）：

\begin{itemize}
  \item \textbf{系統自帶}：iOS [設定] > [私隱與安全性] > [相機]，可關閉「精確位置」；分享嗰陣選「僅當前相片」會\textbf{自動剝離大部分 EXIF}。
  \item \textbf{安卓}：多數品牌相簿 App 喺「分享」嗰陣會提示「係咪移除位置資訊」，勾選即可。
  \item \textbf{專業工具}：ExifTool（電腦端）、Metapho（iOS）、Photo Metadata Remover（安卓）可以徹底清除所有元數據。
\end{itemize}

對於需要\textbf{長期保存}嘅敏感圖片（身份證、銀行卡、戶口本掃描件、護照），建議存入\textbf{加密相簿}：

\begin{itemize}
  \item iPhone：[相片] > 選中相片 > 隱藏（僅 Face ID 驗證後可查看）；
  \item 華為/小米：[相簿] > 長按照片 > 「移入保險櫃」；
  \item 第三方加密 App：\textbf{360 安全相簿}、\textbf{騰訊相冊管家（私隱相簿）}、\textbf{KeepSafe}等，單獨密碼保護，隱蔽性更好。
\end{itemize}
